% Unidad U013. Biografía estructural completa del carácter exponencial.

\label{gfp:b:chap:biografia-completa-e}

\section{Selector orbital finito: existencia y unicidad de la emisión compatible}

Sea \(\mathcal W=\mathbb F_3^6\). Para una emisión
\(U=(u_1,\ldots,u_6)\), la palabra visible es
\[
 \Psi(U)=((-u_1)\bmod3,\ldots,(-u_6)\bmod3).
\]
Sobre \(\mathcal W\), la acción diedral está dada por
\[
\begin{aligned}
 r(u_1,u_2,u_3,u_4,u_5,u_6)
 &=(u_3,u_1,u_2,u_5,u_6,u_4),\\
 s(u_1,u_2,u_3,u_4,u_5,u_6)
 &=(u_4,u_5,u_6,u_1,u_2,u_3).
\end{aligned}
\]
Para \(w\in\mathcal W\), definimos
\[
 n(w)=\#\Psi^{-1}(w),
 \qquad
 M(w)=\sum_{U\in\Psi^{-1}(w)}\operatorname{mult}(U),
\]
y
\[
 \nu(w)=
 \bigl(
 \#\{j:w_j=0\},
 \#\{j:w_j=1\},
 \#\{j:w_j=2\}
 \bigr).
\]

\begin{theorem}[Selección estructural del canal exponencial]
\label{gfp:b:thm:seleccion-estructural-e}
Entre las \(43\) órbitas del catálogo visible, existe una única órbita cuyas
fibras satisfacen simultáneamente:
\[
 n(w)=1,\qquad M(w)=144,\qquad
 \mathfrak D(w)=\{0,3,6\},\qquad
 \nu(w)=(2,3,1).
\]
El calibre
\[
 \operatorname{diag}_c=6,
 \qquad \operatorname{card}=ES
\]
selecciona en ella
\[
 \boxed{w_6^e=201101.}
\]
\end{theorem}

\begin{proof}
Las cuatro condiciones se evalúan sobre las \(243\) palabras del catálogo,
no sobre una muestra. Una única órbita satisface el predicado conjunto. En
esa órbita, una sola emisión posee a la vez diagonal \(6\) y orientación
\(ES\). La enumeración usa únicamente campos del catálogo APP--TRIT y no
consulta una expansión de \(e\).
\end{proof}

\section{Elevación por bloques y compatibilidad de fronteras sucesivas}

La recurrencia \(L_0,L_0,L_1,L_1\) produce
\[
 w_{30}^{e}
 =201101|121221|102011|012222|102011,
\]
y la frontera conjunta aporta
\[
 u_5^e=021222.
\]
El registro de treinta y seis trits queda
\[
 w_{36}^{e}
 =201101|121221|102011|012222|102011|021222.
\]

\section{Firma posicional decimal y cuadrado de compatibilidad}

La lectura inicial de doce cifras es
\[
 \mathbf d_e^{(12)}
 =718281828459
 =7\,\underbrace{1828\,1828}_{\text{cuadrado local}}\,459.
\]
La igualdad central es la concatenación
\[
 1828\mathbin{\|}1828,
\]
no un producto aritmético ni el comienzo supuesto de un período.

\begin{definition}[Cuadrado posicional]
Una palabra \(d_1\cdots d_{12}\) contiene un cuadrado de longitud \(\ell\)
en la posición \(p\) cuando
\[
 d_p\cdots d_{p+\ell-1}
 =d_{p+\ell}\cdots d_{p+2\ell-1}.
\]
La posición y los contextos laterales pertenecen a la firma.
\end{definition}

\begin{table}[htbp]
\centering
\caption{Censo exhaustivo de cuadrados en las lecturas de doce cifras.}
\label{gfp:b:tab:censo-cuadrados-e}
\begin{tabular}{c*{6}{r}}
\toprule
Longitud \(\ell\)&1&2&3&4&5&6\\
\midrule
Ocurrencias&240&21&3&2&0&0\\
Palabras que las contienen&153&19&3&2&0&0\\
\bottomrule
\end{tabular}
\end{table}

La otra palabra con un cuadrado de longitud cuatro es
\[
 575157518472
 =\underbrace{5751\,5751}_{\text{posición inicial }1}\,8472.
\]

\begin{proposition}[Unicidad posicional]
\label{gfp:b:prop:cuadrado-posicional-e}
La lectura del canal \(e\) es la única que posee un cuadrado de longitud
cuatro después de un prefijo de longitud uno y con contextos \(7\) y \(459\).
\end{proposition}

\begin{proof}
El censo deja sólo dos palabras con cuadrados de longitud cuatro. Una lo
inicia en la primera posición; la otra después del prefijo \(7\). La
posición y los contextos separan las dos posibilidades.
\end{proof}

\section{Retorno de fase con cambio de estado y conservación de la memoria recorrida}

En la elevación
\[
 201101|121221|102011|012222|102011
\]
el tercer y el quinto bloques coinciden:
\[
 B_3=B_5=102011.
\]
Sus preestados módulo \(27\) son, sin embargo,
\[
 p_3=(25,11,7,17,8,16),
 \qquad
 p_5=(7,8,4,23,26,7),
\]
y los cocientes posteriores:
\[
 q_3=(6,13,9,2,13,1),
 \qquad
 q_5=(15,0,3,19,23,25).
\]
Se verifican \(p_3\ne p_5\) y \(q_3\ne q_5\).

\begin{figure}[htbp]
\centering
\begin{tikzpicture}[
 node distance=11mm and 7mm,
 every node/.style={font=\small},
 block/.style={draw,rounded corners,minimum width=17mm,minimum height=8mm},
 state/.style={draw,dashed,rounded corners,align=center,inner sep=3pt},
 >=stealth
]
\node[block] (b1) {\(201101\)};
\node[block,right=of b1] (b2) {\(121221\)};
\node[block,right=of b2] (b3) {\(102011\)};
\node[block,right=of b3] (b4) {\(012222\)};
\node[block,right=of b4] (b5) {\(102011\)};
\draw[->] (b1)--(b2);
\draw[->] (b2)--(b3);
\draw[->] (b3)--(b4);
\draw[->] (b4)--(b5);
\node[state,below=of b3] (p3)
 {\(p_3=(25,11,7,17,8,16)\)\\
  \(q_3=(6,13,9,2,13,1)\)};
\node[state,below=of b5] (p5)
 {\(p_5=(7,8,4,23,26,7)\)\\
  \(q_5=(15,0,3,19,23,25)\)};
\draw[->] (p3)--(b3);
\draw[->] (p5)--(b5);
\draw[<->,thick] (b3) to[bend left=28]
 node[above]{misma palabra visible} (b5);
\end{tikzpicture}
\caption{Retorno visible de \(102011\) con renovación del preestado y del
cociente.}
\label{gfp:b:fig:ruptura-memoria-e}
\end{figure}

\begin{theorem}[Retorno visible sin periodicidad del estado]
\label{gfp:b:thm:retorno-visible-no-periodico-e}
\(B_3=B_5\) es un retorno de la proyección visible, no una órbita periódica
del estado enriquecido.
\end{theorem}

\begin{proof}
El retorno del estado completo exigiría igualdad de todas sus coordenadas.
Las desigualdades de preestado y cociente lo excluyen. Tampoco
\(1828|1828\) inicia un período decimal: después del segundo bloque aparece
\(4\), mientras una continuación periódica exigiría \(1\).
\end{proof}

\section{Certificación factorial posterior de la sección exponencial coinductiva}

Sea \(\mathcal A\) un álgebra conmutativa unitaria sobre \(\mathbb Q\), y
sea \(P_t\in\mathcal A[[t]]\) el propagador formal. La concatenación
temporal impone
\[
 P_{s+t}=P_sP_t,
 \qquad
 P_0=1.
\]
Escribimos
\[
 P_t=\sum_{n=0}^{\infty}a_nt^n.
\]
La unidad fija \(a_0=1\), y el transporte elemental normalizado fija
\(a_1=1\).

\begin{theorem}[Recurrencia factorial]
\label{gfp:b:thm:recurrencia-factorial-e}
\[
 (n+1)a_{n+1}=a_n
 \qquad(n\geq0),
\]
y, por tanto,
\[
 a_n=\frac1{n!},
 \qquad
 P_t=\sum_{n=0}^{\infty}\frac{t^n}{n!}.
\]
En \(t=1\),
\[
 \boxed{P_1=e.}
\]
\end{theorem}

\begin{proof}
El coeficiente de \(s^nt\) en \(P_{s+t}\) es
\[
 \binom{n+1}{n}a_{n+1}=(n+1)a_{n+1}.
\]
En \(P_sP_t\) es \(a_na_1=a_n\). La igualdad de series produce la
recurrencia. Una inducción desde \(a_0=1\) da \(a_n=1/n!\).
\end{proof}

\begin{corollary}[Ecuación diferencial]
\[
 \frac{d}{dt}P_t=P_t,
 \qquad P_0=1.
\]
\end{corollary}

\begin{proof}
\[
 P_t'=\sum_{n\geq0}(n+1)a_{n+1}t^n
 =\sum_{n\geq0}a_nt^n=P_t.
\]
\end{proof}

\section{Encaje cofinal de horquillas racionales y convergencia al valor publicado}

Sea
\[
 S_N=\sum_{n=0}^{N}\frac1{n!}.
\]

\begin{theorem}[Cota factorial explícita]
\label{gfp:b:thm:cota-factorial-e}
Para \(N\geq1\),
\begin{equation}
 \boxed{
 S_N<e<
 S_N+\frac{N+2}{(N+1)(N+1)!}.}
 \label{gfp:b:eq:horquilla-factorial-e}
\end{equation}
\end{theorem}

\begin{proof}
El resto es
\[
 R_N=\sum_{k=0}^{\infty}\frac1{(N+1+k)!}>0.
\]
Para \(k\geq0\),
\[
 (N+1+k)!
 \geq(N+1)!(N+2)^k.
\]
Por tanto,
\[
 R_N
 \leq\frac1{(N+1)!}
 \sum_{k=0}^{\infty}\frac1{(N+2)^k}
 =\frac{N+2}{(N+1)(N+1)!}.
\]
La desigualdad es estricta desde \(k=2\).
\end{proof}

Definimos
\[
 J_N^{(e)}=
 \left[
 S_N,\,
 S_N+\frac{N+2}{(N+1)(N+1)!}
 \right]
\]
y, para la parte fraccionaria,
\[
 \widetilde J_N^{(e)}=J_N^{(e)}-2.
\]
La anchura converge a cero. Para cada \(M\), existe \(N\) tal que
\[
 \operatorname{diam}J_N^{(e)}<10^{-M}.
\]
Si los extremos pertenecen a una misma celda decimal, las primeras \(M\)
cifras quedan certificadas mediante aritmética racional.

\Needspace{30\baselineskip}
\section{Prolongación coinductiva de la sección exponencial mediante subcilindros ternarios del Topological Phase Kernel}

En el nivel \(K\), TPK y la relación \(\operatorname{Ext}\) han generado ya
un único subcilindro compatible. Se toma el menor \(N=N_e(K)\) para el cual
la horquilla factorial queda contenida en ese subcilindro. \(N_e(K)\) es un
orden de certificación analítica, no un índice generador ni una regla de
selección de la fibra.

Los cilindros TPK son compatibles bajo truncamiento con independencia del
semigrupo. El teorema semigrupal identifica después su evaluación
arquimediana con \(e-2\); la parte entera recupera \(e\). Si fallan la órbita,
\(\operatorname{Ext}\) o el truncamiento, falla la generación; si fallan la
recurrencia o la cota factorial, falla este certificado analítico, no la
genealogía APP--TRIT--TPK.

\section{Certificado interno de unicidad orbital, recurrencia factorial, convergencia y prolongación cilíndrica}

\begin{remark}[Certificado y control de validez: Falsación del canal exponencial]
La construcción queda refutada si:
\begin{enumerate}
 \item otra órbita satisface el predicado finito completo;
 \item la firma \(7[1828\,1828]459\) no es posicionalmente única;
 \item coinciden los preestados o cocientes de los dos bloques \(102011\);
 \item el semigrupo no fuerza \((n+1)a_{n+1}=a_n\);
 \item la cota \eqref{gfp:b:eq:horquilla-factorial-e} falla;
 \item un nivel no localiza de forma unitaria una extensión compatible;
 \item una expansión tabulada interviene antes del reconocimiento.
\end{enumerate}
\end{remark}
